% !PRINTED 14
\weberpage{14}
%<<CONT
Mannigfaltigkeit angehören, und die denselben beliebig zu wählenden
Nenner haben. Die Addition wird dann so erklärt, dass
\[
 {a\over c}+{b\over c}={a+b\over c}
\]
ist. Diese Regel umfasst als speciellen Fall die Addition der
rationalen Brüche, und um sie allgemein zu rechtfertigen, braucht dann
nur noch gezeigt zu werden, dass, wenn $a:c=a':c'$ und $b:c=b':c'$,
auch $(a+b):c=(a'+b'):c'$ sein muss. Wir beweisen dies, indem wir
zeigen, dass, wenn $a:c=a':c'$ und $(a+b):c>(a'+b'):c'$ ist, auch
$b:c>b':c'$ sein muss. Es sei also, wenn $m$ und $n$ zwei ganze Zahlen
sind,
\[
 {a+b\over c}>{m\over n}>{a'+b'\over c'},
\]
dann ist $n(a+b)>mc$ und also
\[
 {b\over c}>{mc-na\over nc},
 \qquad {mc'-na'\over nc'}>{b'\over c'}.
\]
Andererseits folgt aber leicht aus der Voraussetzung $a:c=a':c'$, dass
auch
\[
 {mc-na\over nc}={mc'-na'\over nc'}
\]
ist, also $b:c>b':c'$, w.~z.~b.~w.

Hiermit ist also nachgewiesen, dass auch die Zahlen, wie wir sie
definirt haben, eine messbare Menge bilden. Sind $a$ und $c$ einer
stetigen Mannigfaltigkeit entnommen, so bilden auch bei feststehendem
$c$ die Verhältnisse $a:c$ eine stetige Menge und es folgt also, da
es überhaupt stetige Mengen giebt, dass auch die Zahlen eine stetige
Menge bilden.

Sind $\alpha,\beta,\gamma,\delta$ jetzt Zahlen, so kann man aus der
Proportion $\alpha:\beta=\gamma:\delta$ eine beliebige der vier Zahlen
durch die drei anderen gegebenen bestimmen. Setzt man $\delta=1$ und
sucht $\alpha$, so erhält man die \gesperrt{Multiplication}
$\alpha=\beta\gamma$, und die Vertauschbarkeit der Factoren ist eine
Folge des Satzes, dass $\alpha:\gamma=\beta:\delta$ aus
$\alpha:\beta=\gamma:\delta$ folgt. Sucht man $\gamma$, so erhält man
die Division; und aus der oben gegebenen Definition der Addition folgt
die Grundformel $\alpha(\beta+\gamma)=\alpha\beta+\alpha\gamma$. Die vier
Grundrechnungsarten sind also in dem Gebiete der Zahlen ausführbar mit
der einzigen Beschränkung, dass bei der Subtraction der Subtrahend
kleiner sein muss als der Minuend.

\begin{center}
\rule{0.18\linewidth}{0.4pt}
\end{center}
